\documentclass[11pt]{article}
\usepackage[a4paper,margin=0.9in]{geometry}
\usepackage[T1]{fontenc}
\usepackage{lmodern}
\usepackage{microtype}
\usepackage{amsmath,amssymb,mathtools}
\usepackage{enumitem}
\usepackage{xcolor}
\usepackage{hyperref}
\usepackage{titlesec}
\usepackage{parskip}

\hypersetup{colorlinks=true, linkcolor=black!70, urlcolor=blue!60!black}

\newcommand{\R}{\mathbb{R}}
\newcommand{\N}{\mathbb{N}}
\newcommand{\dx}{\,dx}
\newcommand{\dt}{\,dt}
\newcommand{\du}{\,du}
\newcommand{\dy}{\,dy}
\newcommand{\dth}{\,d\theta}

\titleformat{\section}{\Large\bfseries\color{blue!40!black}}{}{0em}{}

\newcommand{\wsheader}[3]{%
  \begin{center}\large\textbf{#1 \;--\; #2}\\[2pt]\normalsize\textit{Worksheet #3}\end{center}\medskip
}
\newcommand{\anshdr}[2]{%
  \begin{center}\large\textbf{#1 \;--\; #2}\\[2pt]\normalsize\textit{Answer Key (Worksheets A \& B)}\end{center}\medskip
}
\newenvironment{problems}{\begin{enumerate}[leftmargin=*,label=\textbf{\arabic*.}]}{\end{enumerate}}
\newenvironment{answers}[1]{%
  \medskip\noindent\textbf{Worksheet #1.}
  \begin{enumerate}[leftmargin=*,label=\textbf{\arabic*.}]
}{\end{enumerate}}

\begin{document}

\wsheader{Calc 2}{Convergence Tests}{A}
\begin{problems}
\item Use the integral test to determine whether $\displaystyle \sum_{n=2}^\infty \frac{1}{n \ln n}$ converges.
\item Use the ratio test on $\displaystyle \sum_{n=1}^\infty \frac{n!}{2^n}$.
\item Use the root test on $\displaystyle \sum_{n=1}^\infty \!\Bigl(\tfrac{n}{2n+1}\Bigr)^n$.
\item Use limit comparison on $\displaystyle \sum_{n=1}^\infty \frac{1}{n^2 + n}$.
\item Determine convergence of $\displaystyle \sum_{n=1}^\infty (-1)^n \frac{1}{n}$; is it absolute or conditional?
\end{problems}

\vfill
\newpage
\wsheader{Calc 2}{Convergence Tests}{B}
\begin{problems}
\item Use the integral test on $\displaystyle \sum_{n=1}^\infty \frac{n}{e^n}$.
\item Use the ratio test on $\displaystyle \sum_{n=1}^\infty \frac{2^n}{n!}$.
\item Use the root test on $\displaystyle \sum_{n=1}^\infty \!\Bigl(\tfrac{2n + 1}{3n}\Bigr)^n$.
\item Use direct comparison: $\displaystyle \sum_{n=1}^\infty \frac{\sin^2 n}{n^2}$.
\item Determine convergence of $\displaystyle \sum_{n=1}^\infty (-1)^n \frac{n}{n^2+1}$; is it absolute or conditional?
\end{problems}

\vfill
\newpage
\anshdr{Calc 2}{Convergence Tests}

\begin{answers}{A}
\item $f(x) = 1/(x \ln x)$, with $u = \ln x$: $\int \tfrac{dx}{x \ln x} = \ln(\ln x)$, which $\to \infty$. \textbf{Diverges.}
\item $\tfrac{a_{n+1}}{a_n} = \tfrac{(n+1)!/(2^{n+1})}{n!/2^n} = \tfrac{n+1}{2} \to \infty > 1$. \textbf{Diverges.}
\item $\sqrt[n]{a_n} = \tfrac{n}{2n+1} \to 1/2 < 1$. \textbf{Converges.}
\item Compare with $b_n = 1/n^2$: $a_n/b_n = \tfrac{n^2}{n^2 + n} \to 1$. Since $\sum 1/n^2$ converges, so does $\sum a_n$.
\item Alternating, $b_n = 1/n \downarrow 0$: converges by Leibniz. But $\sum 1/n$ diverges. \textbf{Conditionally convergent.}
\end{answers}

\begin{answers}{B}
\item $\int_1^\infty x e^{-x} \dx = [-xe^{-x} - e^{-x}]_1^\infty = 2/e$. Finite, so the series \textbf{converges}.
\item $\tfrac{a_{n+1}}{a_n} = \tfrac{2}{n+1} \to 0 < 1$. \textbf{Converges.}
\item $\sqrt[n]{a_n} = \tfrac{2n+1}{3n} \to 2/3 < 1$. \textbf{Converges.}
\item $0 \le \tfrac{\sin^2 n}{n^2} \le \tfrac{1}{n^2}$, and $\sum 1/n^2$ converges. \textbf{Converges} by direct comparison.
\item Alternating; $b_n = \tfrac{n}{n^2+1}$ is decreasing for $n \ge 1$ and $\to 0$. So converges by Leibniz. $\sum \tfrac{n}{n^2+1}$ behaves like $\sum 1/n$ (limit comparison): diverges. \textbf{Conditionally convergent.}
\end{answers}

\end{document}
